\section{\label{III_2_B}Jeux de données patrimoniaux} 

The ‘collections as data’ movement seeks to make digital cultural heritage materials available for computational use, encouraging LAM institutions to provide their holdings in machine-actionable formats.Footnote25 For archives, this shift can be valuable, enabling new forms of analysis and engagement. However, as Caswell reminds us, ‘the archive’ in abstract computational discourse is not the same as ‘archives’ as institutions, records, and practices.Footnote26 Treating archives exclusively as data risks flattening their complex provenance, legal status, and social meanings. In this article we therefore adopt a more cautious framing of collections as/for data: we acknowledge that digital surrogates and descriptive metadata can be structured for computational purposes, while insisting that such structuring must remain accountable to archival principles and to the communities whose lives and histories are represented in the records. \footcite{colavizza2026}

La validité d'un modèle d'apprentissage automatique repose sur la qualité et la pertinence des données d'entraînement. Si notre étude a montré la potentielle efficacité d'un pipeline de segmentation puis d'indexation sur un corpus d'images animées spécifique, la généralisation de cette fonctionnalité à un nombre large d'institutions se heurte à l'hétérogénéité des fonds audiovisuels patrimoniaux. 
C'est essentiellement la reconnaissance d'objets fondé sur YoLov8 qui nécessite un jeu de données d'entraînement. L'outil est actuellement entraîné sur des jeux de données d'images génériques, différentes d'images patrimoniales. Les collections conservées au sein des musées et des centres d'archives audiovisuelles se caractérisent en effet par des spécificités visuelles et sémantiques uniques : supports anciens dégradés, esthétiques cinématographiques variées, vocabulaire documentaire hautement spécialisé. Un modèle d'IA entraîné sur un jeu de données restreint risque de produire un surapprentissage (\textit{overfitting}) peu adaptable à la variété des collections gérées par les différents clients de l'éditeur. Pour garantir la robustesse des fonctionnalités automatiques, il s'avère donc indispensable de s'appuyer sur des jeux de données d'entraînement davantage représentatif.
Pour constituer ces corpus d'entraînement à grande échelle sans dépendre des données privées des institutions, la stratégie la plus pertinente réside dans la réutilisation des ressources publiques sous licence libre. 
Dans cette optique, l'entraînement et l'affinage (\textit{fine-tuning}) des modèles de vision par ordinateur peuvent s'appuyer sur d'immenses réservoirs de données ouvertes. Le projet Dataset, créé par l'INA, met notamment à disposition des corpus de données audiovisuelles  "destiné à la mise au point, l’expérimentation et l’évaluation d’outils de recherche et d’analyse de contenus multimédias " \footcite{zotero-item-754}. Ces corpus contiennent des images animées issues de la télévision : journaux télévisés, actualités, captations... Or, ce jeu de données n'entraîne pas nécessairement le modèle sur des images dégradées, en pellicule, en noir et blanc. 
Les Prelinger archives, par exemple, réunissent des archives audiovisuelles américaines de tous types, dénuées de droits d'auteur, y compris des films amateurs ou éducatifs par exemple. \footcite{zotero-item-756}
wikimedia commons 
L'exploitation des API et des corpus sous licence ouverte d'intégrateurs majeurs tels qu'Europeana, le Portail Open Data de la Plateforme Ouverte du Patrimoine (POP) du Ministère de la Culture, ou les fonds numérisés de Gallica (BnF).

Double entraînement : données + règles métier (FRBR) : pour le NLP : LLM

L'entraînement des modèles reposant sur un fonds unique d'un client donné, comme nous l'avons exposé en seconde partie, poserait plusieurs difficultés pour la généralisation de la fonctionnalité. D'une part, les données d'une seule institution ne sont pas suffisantes ; d'autre part, cela exige un accord contractuel explicite de cession de droits d'usage sur les données. La création d'un modèle mutuel suppose une gouvernance transparente, où chaque institution accepte d'alimenter un réservoir partagé de \textit{Ground Truth} dans une logique de mutualisation des ressources.

Enfin, cette réflexion sur les jeux de données d'entraînement résonne directement avec les autres briques technologiques en cours de spécification au sein de l'entreprise, notamment la recherche en langage naturel et la recherche visuelle par similarité. La construction d'un jeu de données patrimonial universel — associant visuels sous licence libre, descripteurs FRBR et thésaurus normés — profite à l'ensemble du système d'information :
\begin{enumerate}
	\item Les métadonnées textuelles ouvertes renforcent le \textit{mapping} sémantique de la recherche en langage naturel.
	\item Les paires images-concepts réutilisables affinent les réseaux de neurones dédiés à la recherche par l'image et à la segmentation vidéo.
\end{enumerate}

En plaçant la réutilisation des données ouvertes et le principe des \textit{Collections as Data} au cœur de sa R\&D, Skinsoft garantit la soutenabilité technique de ses fonctionnalités automatiques tout en respectant scrupuleusement le cadre légal et éthique de ses clients patrimoniaux.



choix d'appui sur des données produites directement par des institutions culturelles. utiliser exclusivement des sources ouvertes. pour les images : entrainement sur des images disponibles sous licence ouverte : european, POP + textes de lois, méthodologie des gestion des collections, thesaurus, normes de catalogage...R
Il faut l'accord client et les données d'un seul client ne sont pas suffisantes.